\documentclass[../main.tex]{subfiles}

\begin{document}

\chapter{微分与积分}\label{chapter:ch06-微分与积分}

与 Riemann 积分有关的积分与微分的计算的基本定理是微积分学的基石. 本章我们对 Lebesgue 积分阐述这两个定理. 对于闭有界区间 $[a,b]$ 上的函数 $f$, 何时有
\begin{align}
\int_{a}^{b}f'=f(b)-f(a)?\label{align:ch06-积分基本定理}
\end{align}
假设 $f$ 是连续的. 延拓 $f$ 使得在 $(b,b+1)$ 上取值 $f(b)$, 对于 $0<h\leqslant 1$, $[a,b]$ 上的均差函数 $\operatorname{Diff}_{h}f$ 和平均值函数 $\operatorname{Avg}_{h}f$ 定义为
\begin{align*}
\text{对所有在}[a,b]\text{内的 }x,\quad\operatorname{Diff}_{h}f(x)=\frac{f(x+h)-f(x)}{h}\quad\text{与}\quad\operatorname{Avg}_{h}f(x)=\frac{1}{h}\int_{x}^{x+h}f(t)\,dt.
\end{align*}
通过变量的代换和一些项的抵消, 对于 Riemann 积分给出了\eqref{align:ch06-积分基本定理}的离散形式:
\begin{align*}
\int_{a}^{b}\operatorname{Diff}_{h}f=\operatorname{Avg}_{h}f(b)-\operatorname{Avg}_{h}f(a).
\end{align*}
当 $h\to 0^{+}$ 时, 右边的极限等于 $f(b)-f(a)$. 我们证明 Henri Lebesgue 的一个震撼性的定理, 它告诉我们 $(a,b)$ 上的单调函数几乎处处有有限导数. 接着定义对于一个函数绝对连续意味着什么, 且证明: 若 $f$ 绝对连续, 则 $f$ 是两个单调函数的差且均差族 $\{\operatorname{Diff}_{h}f\}_{0<h\leqslant 1}$ 是一致可积的. 因此, 根据 Vitali 收敛定理, 通过在离散形式取 $h\to 0^{+}$ 的极限得到 $f$ 绝对连续, 所以\eqref{align:ch06-积分基本定理}成立. 若 $f$ 单调且\eqref{align:ch06-积分基本定理}成立, 我们证明 $f$ 必须是绝对连续的. 从积分形式的基本定理\eqref{align:ch06-积分基本定理}, 我们得到微分形式的基本定理, 即若 $f$ 在 $[a,b]$ 上是 Lebesgue 可积的, 则对 $[a,b]$ 中的几乎所有 $x$,
\begin{align}
\frac{d}{dx}\left[\int_{a}^{x}f\right]=f(x).\label{align:ch06-微分基本定理}
\end{align}

\subfile{Chapter-06/section-01.tex}

\subfile{Chapter-06/section-02.tex}

\subfile{Chapter-06/section-03.tex}

\subfile{Chapter-06/section-04.tex}

\subfile{Chapter-06/section-05.tex}

\subfile{Chapter-06/section-06.tex}

\end{document}